% Physical unit 019; accepted Session 02.
\GordanUnitStart{019}{printed page 14}
\GordanTopFolio{14}
\[
=(10)+{(11)\cos^2l\over \cos l_0\sqrt{1-e^2\cos^2l}}
\left\{{\alpha\sin l_0\cos\varphi+\beta{du\over d\varphi}
\over \alpha\sin l_0\sin\varphi+\beta u}
-{ {dv\over d\varphi}\over v}\right\}.
\]

\begin{center}\S\ 8.\end{center}

Since on the right side of the equation the quantities
\[
 (01)\quad (11)\quad \cos^2l\quad \tg l_0\quad {1\over \sqrt{1-e^2\cos^2l}}\quad
 \cos\varphi\quad {du\over d\varphi}\quad {dv\over d\varphi}\quad {1\over v}
\]
remain perpetually finite, the function $N$ can become infinite only in this way, namely if
\[
 \alpha\sin l_0\sin\varphi+\beta u=0
\]
holds.

It follows from this that the integral
\[
 \int_{\psi_1}^{\psi_2}d\varphi\{(11)\delta l'\delta l'+2(01)\delta l\delta l'+(00)\delta l\delta l\}
\]
is positive so long as the constants $\alpha$ and $\beta$ can be determined in such a way that
\[
 \alpha\sin l_0\sin\varphi+\beta u
\]
does not vanish; thus so long as
\[
 {u\over \sin l_0\sin\varphi}
\]
does not assume all numbers lying between $-\infty$ and $+\infty$. For if
\[
 {u\over \sin l_0\sin\varphi}
\]
does not assume some value ``$a$'' in this interval, then, putting
\[
 \alpha+\beta a=0,
\]
there is never
\[
 \alpha\sin l_0\sin\varphi+\beta u=0.
\]

\begin{center}\S\ 9.\end{center}

Since
\[
 {d\over d\varphi}\left\{ {u\over \sin l_0\sin\varphi}\right\}
 =-{1-\sin^2l_0\cos^2\varphi\over \sin l_0\cos l_0\sin^2\varphi}
\]
is always negative, and since, on putting $\psi=\psi_1$,
\[
 {u\over \sin l_0\sin\varphi}=0,
\]
as $\psi$ increases,
\[
 {u\over \sin l_0\sin\varphi}
\]
assumes in succession all values from $0$ down to $-\infty$ and, after passing back from $-\infty$ to $+\infty$, from $\infty$ to $0$.
